\chapter{Method, scope, and evidence classification}
\label{method-scope-evidence}

This work is an evidence-grounded technical monograph rather than a comprehensive systematic review. It asks which claims about the reliability of coding-agent systems can be supported, which operational practices follow from those claims, and which observations a team must reproduce locally before adopting a practice. The unit of analysis is the system around an agent: its evaluation harness, execution and recovery machinery, retrieval and context path, permission boundary, review interface, and allocation policy.

\section{Evidence collection and scope}

The scholarly corpus contains 118 sources organized into seven topic-specific threads: benchmark validity, failure taxonomies, evaluation statistics, human oversight and accountability, context and retrieval, durable execution, and scheduling with repository-scale scoping. The searches were consolidated on July 26, 2026, using the SciX scholarly index, which mirrors arXiv and NASA ADS metadata. Six threads contain 17 sources each; a seventh contains 16 sources added after an adversarial citation audit identified relevant scheduling and repository-scoping work that had entered the draft through other syntheses. Metadata and identifiers were checked against the scholarly index.

The external operational corpus adds 91 practitioner records drawn from engineering reports, incident accounts, vendor and independent analyses, conference material, and community field reports. A separate benchmark catalog records 29 coding-agent and adjacent evaluation suites. Seventeen project dossiers describe author-operated systems and experiments. These sources were selected to cover the system boundaries above, not to estimate a prevalence rate for the entire literature. The resulting corpus should therefore be read as a structured, current evidence synthesis, not as an exhaustive or publication-bias-corrected census.

\section{Evidence classification}

Each admitted item carries two independent labels. The \emph{source class} records where the claim comes from: research literature, an explorer synthesis of a topic-scoped source set, or a practitioner account. The \emph{strength grade} records what the item can support. A \textbf{strong} item directly supports the stated practice through a controlled result, peer-reviewed methodological result, or independently recomputable artifact. A \textbf{directional} item supports the direction of a claim while leaving its magnitude or transfer uncertain. An \textbf{anecdotal} item establishes a reported case, not its frequency. A \textbf{null-result} item records a test that did not support the proposed effect. Contested items retain their counterevidence rather than being resolved by averaging incompatible claims.

Source class and strength are deliberately separate. A research paper may be too narrow to support a general operational prescription, while a recomputable project artifact may strongly support a local factual claim. Every taught practice reports its evidence at the start of the relevant treatment and lists the supporting records in a chapter-level ``Sources and evidence'' section.

A separate evidence audit rechecked every explorer-synthesis item initially graded strong against its underlying source. The grade was retained only when the source reported a controlled result that supported the practice's own claim; otherwise the item was downgraded or its prose was narrowed. This audit prevents agreement among summaries from being mistaken for independent experimental support.

\section{Practice derivation and selection}

The practice catalog was derived in two independent passes. In the first, eight specialist analyses examined separate slices of the corpus: fleet operations, durable execution, evaluation practice, measurement methodology, failure diagnostics, context and structure, attention and accountability, and existing-draft material. Their candidate practices were merged and adversarially checked line by line against sources; 43 evidence corrections were recorded during that stage. In the second pass, a whole-corpus analysis derived practices without seeing the specialist outputs. Agreement between the two differently structured passes is used as a derivation-confidence signal, not as a substitute for evidence strength.

A candidate practice was admitted only when it named an action, a mechanism, and a boundary at which the action could stop working. This process produced 192 practices grouped into 14 mechanism clusters. Three independent selection lenses then proposed which practices required full treatment by asking whether a practice changes an engineering decision, whether omitting it loses mechanism coverage, and whether the available case can teach the mechanism without exceeding its evidence. Their agreements and disagreements determined the 55 practices developed in the main chapters; the remainder stay in the companion catalog as bounded alternatives or thin-support research prompts.

\section{Contributions}

The monograph makes five contributions:

\begin{enumerate}
  \item an audited evidence map that separates source provenance from evidentiary strength and preserves null, anecdotal, and contested results;
  \item a catalog of 192 bounded engineering practices, with 55 practices developed in depth and the remainder retained for scoped follow-up;
  \item a dependency chain across measurement, grading, containment and recovery, retrieval and context, human oversight, and allocation, showing what each operational layer permits the next layer to claim;
  \item scoped measurements and failure analyses from author-operated coding-agent and retrieval systems, reported with their local conditions and used to expose mechanisms rather than to estimate universal effects; and
  \item runnable protocols for repeated evaluation, paired comparison, oracle and grader validation, fault injection, trace analysis, retrieval diagnosis, escalation design, topology comparison, and scheduling replay.
\end{enumerate}

\section{Treatment of author-system cases}

Cases from the author's systems are illustrations. They reveal state transitions, permission boundaries, measurement defects, or recovery failures and can motivate a locally reproducible test. They are not counted as independent external evidence for a general recommendation. Where an author-system result is reported, the claim is restricted to the named workload, task set, scorer, execution window, and available artifact. Chapter-level evidence sections identify when a case corroborates a mechanism but does not admit the practice into the evidence base.

\section{Limitations}

The corpus is current through July 2026 but is not exhaustive. Search threads were selected around the monograph's engineering questions, and the practitioner corpus is especially vulnerable to selective reporting. The evidence is uneven: measurement and retrieval have substantially more controlled support than containment, recovery, and fleet allocation. Several results concern language, office, or general agents rather than repository-scale coding agents, and their transfer is labeled rather than assumed.

Fast-moving model and harness capabilities also limit the life of numerical results. Context lengths, grader error rates, localization accuracy, model costs, and routing frontiers are snapshots rather than constants. The protocols are intended to make those quantities locally measurable. They do not supply universal thresholds where the corpus has none.

Finally, the derivation process used multiple model-assisted analyses to inspect and compare the corpus. Independence here describes separation of inputs and visibility between derivation passes, not statistical independence or human inter-rater reliability. All retained claims remain traceable to the listed human-authored sources or author-owned artifacts, and the evidence grades describe those sources rather than the confidence of the model-assisted derivation process.
